% !TEX root = ../main.tex
\chapter{Annexes Techniques et Documentation}
\label{annexe:documentation}

\section{A.1 — Code Source Commenté du Firmware ESP32}

\subsection{A.1.1 — Fichier Principal : \texttt{main.cpp}}

\begin{tcolorbox}[colback=gray!5, colframe=macouleur, title=\texttt{main.cpp} -- Point d'entrée du firmware ESP32]
\begin{verbatim}
/**
 * @file main.cpp
 * @brief Firmware CNC 5-axes pour ESP32 (configuration Trunnion Table)
 * @author Étudiant PFE - ENI-ABT, Génie Mécanique et Énergie
 * @date 2025
 * @hardware ESP32-WROOM-32U, DM556 x5, NEMA17 x5, MeanWell 36V/350W
 *
 * ARCHITECTURE DUAL-CORE:
 *   - Core 0: Réseau Wi-Fi (UDP), parsing JSON, cinématique inverse
 *   - Core 1: Hardware Timer ISR (Step/Dir), Ring Buffer, Planner
 */

#include <Arduino.h>
#include <WiFi.h>
#include <AsyncUDP.h>
#include "GCodeParser.h"       // Parser G-Code ISO 6983
#include "Planner5Axes.h"      // Planificateur Look-Ahead
#include "InvKinematics.h"     // Cinématique Inverse Trunnion
#include "StepGenerator.h"     // Générateur d'impulsions Step/Dir
#include "SafetyMonitor.h"     // Gestion E-STOP et fins de course

// ---- Configuration Réseau Wi-Fi ----
const char* WIFI_SSID = "CNC5AXES_AP";
const char* WIFI_PASS = "cnc12345";    // Réseau AP propre à la machine
const uint16_t UDP_PORT = 2222;

// ---- Instances Globales ----
AsyncUDP udp;
GCodeParser  parser;
Planner5Axes planner;
InvKinematics ik;
StepGenerator stepGen;
SafetyMonitor safety;

// Tâche réseau sur Core 0 (basse priorité)
void taskNetwork(void* pvParameters) {
    WiFi.softAP(WIFI_SSID, WIFI_PASS);
    udp.listen(UDP_PORT);
    
    udp.onPacket([](AsyncUDPPacket packet) {
        // Extraire la charge utile JSON
        String payload((char*)packet.data(), packet.length());
        
        // Valider le CRC-16 du paquet
        if (!validateCRC(payload)) return;
        
        // Parser le bloc G-Code
        GCodeBlock block = parser.parseBlock(payload);
        
        // Calcul cinématique inverse (TCP compensé)
        MotorPositions motors = ik.solve(
            block.x, block.y, block.z,
            block.a, block.b);
        
        // Insérer dans le Ring Buffer du Planner
        planner.bufferBlock(motors, block.feedRate);
        
        // Envoyer ACK avec position courante
        String ack = buildACKJson(ik.getCurrentPosition());
        packet.print(ack);
    });
    
    // Boucle infinie (la pile réseau gère l'événementiel)
    while(true) vTaskDelay(100 / portTICK_PERIOD_MS);
}

// Tâche temps-réel sur Core 1 (priorité maximale)
void taskRealTime(void* pvParameters) {
    // Configurer les Hardware Timers pour Step Generation
    stepGen.init();
    safety.init();
    
    while(true) {
        // Vérifier les fins de course et E-STOP
        if (safety.isEmergencyStop()) {
            stepGen.disableAllAxes();
            planner.clearBuffer();
            // Attendre l'acquittement de l'opérateur
            while(safety.isEmergencyStop()) vTaskDelay(1);
        }
        
        // Exécuter le bloc suivant s'il est disponible
        if (planner.hasBlock()) {
            MotionBlock block = planner.popBlock();
            stepGen.executeBlock(&block);  // Bloquant
        }
        
        vTaskDelay(1 / portTICK_PERIOD_MS); // Yield 1ms
    }
}

void setup() {
    Serial.begin(115200);
    
    // Démarrer les deux tâches sur leurs cœurs assignés
    xTaskCreatePinnedToCore(taskNetwork,  "NetTask",  8192, NULL, 1, NULL, 0);
    xTaskCreatePinnedToCore(taskRealTime, "RTTask",   4096, NULL, 5, NULL, 1);
    
    // Déléguer immédiatement à FreeRTOS Scheduler
}

void loop() { vTaskDelete(NULL); } // Loop non utilisée
\end{verbatim}
\end{tcolorbox}

\subsection{A.1.2 — Algorithme de Cinématique Inverse : \texttt{InvKinematics.cpp}}

\begin{tcolorbox}[colback=gray!5, colframe=macouleur, title=\texttt{InvKinematics.cpp} -- Calcul des compensations TCP Trunnion]
\begin{verbatim}
/**
 * @brief Résout la cinématique inverse de la table Trunnion.
 * Transforme les coordonnées pièce (x_w, y_w, z_w, A°, B°)
 * en positions machine corrigées (x_m, y_m, z_m).
 *
 * ÉQUATIONS MATHÉMATIQUES (Trunnion A+B):
 *   R_plateau = rayon du plateau = 100 mm
 *   H_berceau = hauteur du pivot sous la table = 80 mm
 *
 *   dZ = R * (1 - cos(A_rad))       // Décalage TCP en Z
 *   dY = R * sin(A_rad)              // Décalage TCP en Y
 *
 *   x_m = x_w   - dY * (-sin(B_rad)) // Projection X par rotation B
 *   y_m = y_w   - dY * cos(B_rad)    // Projection Y
 *   z_m = z_w   - dZ                 // Correction Z
 */
MotorPositions InvKinematics::solve(float x_w, float y_w, float z_w,
                                     float A_deg, float B_deg) {
    const float R = PLATEAU_RADIUS_MM;  // 100.0f mm
    
    // Conversion Degrés -> Radians
    float A_rad = A_deg * DEG_TO_RAD;
    float B_rad = B_deg * DEG_TO_RAD;
    
    // Calcul des offsets TCP (Table Trunnion A+B)
    float dZ_tcp = R * (1.0f - cosf(A_rad));
    float dY_tcp = R * sinf(A_rad);
    
    // Application de la rotation B sur les offsets (matrice 2x2)
    float dx_machine = dY_tcp * (-sinf(B_rad));
    float dy_machine = dY_tcp * cosf(B_rad);
    
    // Positions machine finales
    MotorPositions motors;
    motors.x = x_w - dx_machine;
    motors.y = y_w - dy_machine;
    motors.z = z_w - dZ_tcp;
    motors.a = A_deg;  // Angles moteurs A et B
    motors.b = B_deg;  // = angulations demandées
    
    return motors;
}
\end{verbatim}
\end{tcolorbox}

\subsection{A.1.3 — Interface Utilisateur Flutter : \texttt{dashboard\_page.dart}}

\begin{tcolorbox}[colback=gray!5, colframe=macouleur, title=\texttt{dashboard\_page.dart} -- Écran principal IHM Flutter]
\begin{verbatim}
/// Page principale de l'IHM - Tableau de bord de la CNC 5 Axes
/// Architecture : Riverpod v2 + Flutter 3.x
/// Widget Stateless consommant les états depuis les Providers
class DashboardPage extends ConsumerWidget {
  const DashboardPage({super.key});

  @override
  Widget build(BuildContext context, WidgetRef ref) {
    // Écoute réactive de l'état complet de la machine
    final machineState = ref.watch(machineStateProvider);
    final coordinates = ref.watch(coordinatesProvider);
    final connectionStatus = ref.watch(connectionStatusProvider);

    return Scaffold(
      backgroundColor: const Color(0xFF0A0F1E),
      body: Row(
        children: [
          // === PANNEAU GAUCHE : Coordonnées & Status ===
          SizedBox(
            width: 280,
            child: Column(
              children: [
                // En-tête Machine
                MachineHeaderWidget(
                  machineName: 'CNC-5AXES / ENI-ABT',
                  status: connectionStatus,
                ),
                // Affichage des 5 axes (monospace, grand)
                for (final axis in ['X', 'Y', 'Z', 'A', 'B'])
                  AxisCoordinateCard(
                    axisLabel: axis,
                    value: coordinates.getAxis(axis),
                    isRotary: axis == 'A' || axis == 'B',
                  ),
              ],
            ),
          ),
          
          // === PANNEAU CENTRE : Visualisateur 3D G-Code ===
          Expanded(
            child: GCodeViewer3D(
              gcodeBlocks: ref.watch(gcodeQueueProvider),
              currentLineIndex: machineState.currentLine,
            ),
          ),
          
          // === PANNEAU DROITE : Contrôles ===
          SizedBox(
            width: 300,
            child: Column(
              children: [
                // Boutons Action Principaux
                CycleControlButtons(
                  onStart: () => ref.read(
                    cycleControlProvider.notifier).startCycle(),
                  onPause: () => ref.read(
                    cycleControlProvider.notifier).feedHold(),
                  onEStop: () => ref.read(
                    safetyProvider.notifier).triggerEStop(),
                ),
                // Contrôle Jog Manuel
                const JogPadWidget(),
              ],
            ),
          ),
        ],
      ),
    );
  }
}
\end{verbatim}
\end{tcolorbox}

% ============================================
\section{A.2 — Fiches Techniques des Composants Principaux}

\subsection{A.2.1 — Extrait de la Fiche Technique ESP32-WROOM-32U (Espressif Systems)}

\begin{table}[htbp]
\centering
\renewcommand{\arraystretch}{1.3}
\begin{tabular}{|l|l|}
\hline
\rowcolor{macouleur!20}
\textbf{Paramètre} & \textbf{Valeur} \\
\hline
Processeur & Xtensa LX6 Dual-Core 32-bit \\
\hline
Fréquence CPU & 80 / 160 / 240 MHz (configurable) \\
\hline
SRAM interne & 520 KB \\
\hline
Flash embarquée & 4 MB (SPI Flash) \\
\hline
Wi-Fi & 802.11 b/g/n 2.4 GHz, 150 Mbps \\
\hline
Bluetooth & BT Classic 4.2 + BLE \\
\hline
GPIO disponibles & 34 broches (28 configurables) \\
\hline
Hardware Timers & 4 × 64 bits haute résolution \\
\hline
ADC & 18 canaux 12 bits (0--3.6V) \\
\hline
DAC & 2 canaux 8 bits \\
\hline
Tension d'alimentation & $3{,}0$ -- $3{,}6$ V (Typique $3{,}3$ V) \\
\hline
Courant de veille (Deep Sleep) & $10\,\mu A$ \\
\hline
Courant opération Wi-Fi (Pic Tx) & $500\,mA$ (Burst), $80\,mA$ nominal \\
\hline
Température opération & $-40°C$ -- $+85°C$ (Gamme Industrielle) \\
\hline
Dimensions module & $18 \times 25{,}5 \times 3{,}1$ mm \\
\hline
Certification & CE, FCC, TELEC, IC \\
\hline
\end{tabular}
\caption{Caractéristiques techniques clés de l'ESP32-WROOM-32U (Source : Espressif DS ESP32 Rev 3.3)}
\label{tab:esp32_specs}
\end{table}

\subsection{A.2.2 — Extrait de la Fiche Technique Driver DM556 (Leadshine)}

\begin{table}[htbp]
\centering
\renewcommand{\arraystretch}{1.3}
\begin{tabular}{|l|l|}
\hline
\rowcolor{macouleur!20}
\textbf{Paramètre} & \textbf{Valeur} \\
\hline
Courant moteur (pic) & 1{,}0 -- 5{,}6 A (réglable par DIP) \\
\hline
Tension d'alimentation & 18 -- 50 V DC \\
\hline
Micro-stepping & 400 / 800 / 1~600 / 3~200 / 6~400 / 12~800 / 25~600 pas/tour \\
\hline
Fréquence Step max & 200 kHz \\
\hline
Temps de réponse Step & $< 1\,\mu s$ (Isolation optique) \\
\hline
Niveau logique entrée & $3{,}5$ -- $5{,}5$ V (TTL 5V requis) \\
\hline
Courant d'entrée ($I_f$) & $10$ -- $14$ mA \\
\hline
Protection thermique & Oui ($75°C$ coupure auto) \\
\hline
Protection court-circuit & Oui \\
\hline
Protection surtension & Oui (Moteur back-EMF) \\
\hline
Rendement & $\approx 95\%$ \\
\hline
Dimensions & $118 \times 75.5 \times 34$ mm \\
\hline
\end{tabular}
\caption{Caractéristiques techniques clés du Driver DM556 (Source : Leadshine DS 2018)}
\label{tab:dm556_specs}
\end{table}

\subsection{A.2.3 — Caractéristiques du Moteur NEMA 17HS4401}

\begin{table}[htbp]
\centering
\renewcommand{\arraystretch}{1.3}
\begin{tabular}{|l|l|}
\hline
\rowcolor{macouleur!20}
\textbf{Paramètre} & \textbf{Valeur} \\
\hline
Pas angulaire (plein pas) & $1{,}8°$ (200 pas/tour) \\
\hline
Holding Torque & $0{,}45$ Nm \\
\hline
Courant de phase (nominal) & $1{,}7$ A \\
\hline
Résistance de phase & $1{,}5\,\Omega$ \\
\hline
Inductance de phase & $3{,}6$ mH \\
\hline
Inertie rotor & $68\,g.cm^2$ \\
\hline
Longueur (hors arbre) & 40 mm \\
\hline
Poids & 0{,}34 kg \\
\hline
Taille boîtier & $42{,}3 \times 42{,}3$ mm (NEMA 17) \\
\hline
Interface arbre & $\varnothing\, 5$ mm (simple méplat) \\
\hline
\end{tabular}
\caption{Caractéristiques techniques du moteur pas-à-pas NEMA 17HS4401}
\label{tab:nema17_specs}
\end{table}

% ============================================
\section{A.3 — Données Brutes des Tests de Métrologie (Axe Y -- 20 Mesures)}

\begin{table}[htbp]
\centering
\renewcommand{\arraystretch}{1.3}
\begin{tabular}{|c|c|c||c|c|c|}
\hline
\rowcolor{macouleur!20}
\textbf{N° Mesure} & \textbf{Valeur ($\mu m$)} & \textbf{Écart ($\mu m$)} & \textbf{N° Mesure} & \textbf{Valeur ($\mu m$)} & \textbf{Écart ($\mu m$)} \\
\hline
1 & $-26$ & $+3$ & 11 & $-28$ & $+1$ \\
2 & $-31$ & $-2$ & 12 & $-32$ & $-3$ \\
3 & $-29$ & $0$ & 13 & $-27$ & $+2$ \\
4 & $-28$ & $+1$ & 14 & $-30$ & $-1$ \\
5 & $-33$ & $-4$ & 15 & $-29$ & $0$ \\
6 & $-27$ & $+2$ & 16 & $-27$ & $+2$ \\
7 & $-29$ & $0$ & 17 & $-32$ & $-3$ \\
8 & $-30$ & $-1$ & 18 & $-28$ & $+1$ \\
9 & $-26$ & $+3$ & 19 & $-31$ & $-2$ \\
10 & $-29$ & $0$ & 20 & $-29$ & $0$ \\
\hline
\multicolumn{2}{|c|}{\textbf{Moyenne}} & $\boldsymbol{\bar{x} = -29}$ & \multicolumn{2}{c|}{\textbf{Déviation Std.}} & $\boldsymbol{\sigma = 1{,}9\,\mu m}$ \\
\hline
\end{tabular}
\caption{Données brutes des 20 mesures métrologie de l'axe Y (déplacement consigne = 50 mm)}
\label{tab:raw_data_y}
\end{table}

% ============================================
\section{A.4 — Manuel Opérateur (Guide de Démarrage Rapide)}

\subsection{A.4.1 — Procédure de Démarrage}
\begin{enumerate}
    \item \textbf{Vérifications avant mise sous tension :}
    \begin{itemize}
        \item S'assurer que la zone de travail est libre d'obstacles.
        \item Vérifier que le bouton E-STOP (champignon rouge) est \textbf{enfoncé}.
        \item Contrôler visuellement les connexions des câbles moteurs.
    \end{itemize}
    \item \textbf{Mise sous tension :} Activer l'interrupteur général du coffret. L'afficheur OLED affiche «BOOTING...» puis «IDLE — WIFI OK».
    \item \textbf{Connexion Wi-Fi :} Connecter l'ordinateur/tablette au réseau \textbf{CNC5AXES\_AP} (Mot de passe: \texttt{cnc12345}).
    \item \textbf{Lancement de l'IHM Flutter :} Ouvrir l'application. L'indicateur de connexion passe au vert.
    \item \textbf{Déverrouillage E-STOP :} Tourner le champignon rouge d'un quart de tour (sens anti-horaire). Cliquer sur «Déverrouiller (\$X)» dans l'IHM.
    \item \textbf{Homing (Référencement) :} Cliquer sur «Homing (\$H)» dans l'IHM. La machine va automatiquement referencer les 5 axes sur leurs fins de course. La LED verte «CONNECTÉ» s'allume.
    \item \textbf{Chargement du fichier G-Code :} Utiliser l'onglet «Éditeur» pour charger le fichier \texttt{.nc}. Vérifier le résultat de la validation.
    \item \textbf{Lancement de l'usinage :} Appuyer sur «LANCER L'USINAGE» (bouton vert). Surveiller l'avancement dans le visualisateur 3D.
\end{enumerate}

\subsection{A.4.2 — Procédure d'Arrêt d'Urgence}
En cas de comportement anormal (vibrations anormales, odeur de brûlé, collision imminente) :
\begin{enumerate}
    \item Appuyer immédiatement sur le champignon rouge E-STOP.
    \item Cela coupe l'alimentation 36V et désactive tous les moteurs.
    \item Ne pas relancer la machine avant d'avoir identifié et corrigé la cause.
    \item Un redémarrage complet (séquence depuis l'étape 1) est requis.
\end{enumerate}

\subsection{A.4.3 — Planning de Maintenance Préventive}

\begin{table}[htbp]
\centering
\renewcommand{\arraystretch}{1.3}
\begin{tabular}{|c|p{5cm}|l|p{4cm}|}
\hline
\rowcolor{macouleur!20}
\textbf{Fréquence} & \textbf{Opération} & \textbf{Composant} & \textbf{Produit/Outil} \\
\hline
Avant chaque utilisation & Inspection visuelle des câbles & Câbles moteurs & Visuel \\
\hline
Avant chaque utilisation & Serrage des guides & Vis de serrage HGR15 & Clé Allen 4mm \\
\hline
Hebdomadaire & Nettoyage des copeaux & Fraise, rails, plateau & Soufflette + chiffon \\
\hline
Mensuel & Graissage des rails HGR15 & Patins à billes & Graisse lithium EP2 \\
\hline
Mensuel & Graissage des écrous vis à billes & SFU1605 x3 & Huile machine SAE 10 \\
\hline
Trimestriel & Test E-STOP (Fonctionnel) & Contacteur KM1 & Manuel \\
\hline
Semestriel & Vérification des courroies (tension) & Entraînements A+B & Tensiomètre \\
\hline
Annuel & Recalibration complète (métrologie) & 5 axes & Comparateur $1\,\mu m$ \\
\hline
Aux 500 h & Remplacement courroies GT2 & Entraînements A+B & Courroie GT2-6mm neuve \\
\hline
Aux 1~000 h & Remplacement joints OLED & Afficheur OLED & SSD1306 neuf \\
\hline
\end{tabular}
\caption{Planning de maintenance préventive de la CNC 5 axes}
\label{tab:maintenance}
\end{table}
